\documentclass[11pt]{article}

\usepackage[a4paper,margin=1in]{geometry}
\usepackage[T2A]{fontenc}
\usepackage[utf8]{inputenc}
\usepackage[russian]{babel}
\usepackage{amsmath,amssymb,amsthm}
\usepackage{booktabs}
\usepackage{hyperref}
\usepackage{microtype}

\newtheorem{theorem}{Теорема}
\newtheorem{proposition}[theorem]{Предложение}
\newtheorem{corollary}[theorem]{Следствие}
\newtheorem{lemma}[theorem]{Лемма}
\theoremstyle{remark}
\newtheorem{remark}[theorem]{Замечание}

\title{diff-gpt: авторегрессионное прогнозирование временных рядов \\
       с помощью токенизации производных и $\Sigma\Delta$-квантования}
\author{Михаил Борисов}
\date{\today}

\begin{document}
\maketitle

\begin{abstract}
Мы изучаем авторегрессионное прогнозирование непрерывных временных рядов
decoder-only трансформером, обучаемым на токенизации на основе производных.
Энкодер дифференцирует вход, масштабирует его по наблюдаемому диапазону и
квантует $\Sigma\Delta$-модулятором первого порядка; декодер точно обращает
эту процедуру с точностью до слагаемого порядка ширины словаря. Мы доказываем
две оценки на уровне токенизации: достижимая кросс-энтропийная ошибка
растёт как $\log V$, а ошибка реконструкции убывает как $1/V$ --- обе
независимо от модели. Дополнительно вводится порядковая кросс-энтропия
с гауссовской мягкой меткой, превращающая классификацию по бинам в
регрессию с учётом расстояния между бинами. На бенчмарке M4 Hourly
порядковое сглаживание снижает sMAPE с $18{,}58$ до $17{,}49$ (6\%
относительного улучшения) без изменения архитектуры и расписания обучения.
\end{abstract}

\section{Введение}

Большинство современных моделей прогнозирования временных рядов ---
это регрессионные модели со специализированными индуктивными смещениями
(patching, inverted attention, channel mixing). Они рассматривают выход
как непрерывный вектор и минимизируют $\ell_2$ или Huber-потерю. Большие
языковые модели, напротив, суть decoder-only трансформеры, обучаемые
кросс-энтропией над дискретным словарём. Вопрос, которым мы задаёмся:
конкурентоспособен ли чистый LM-рецепт --- фиксированная архитектура,
фиксированная функция потерь, авторегрессионная генерация --- на
стандартных бенчмарках прогнозирования, если непрерывный вход
предварительно отображён в точное дискретное представление.

Ключевым выбором является токенизация. Мы описываем энкодер на основе
производной с $\Sigma\Delta$-квантованием первого порядка, обладающий
двумя важными свойствами: (а) кросс-энтропийная ошибка на достижимом
пределе масштабируется как $\log V$, поэтому размер словаря имеет
предсказуемую информационную цену, и (б) ошибка реконструкции пары
энкодер--декодер равномерно ограничена по длине последовательности
величиной $\Delta / 2 = D / (V - 2)$, где $D$ --- максимальная
наблюдаемая первая разность. Обе оценки представляют собой
токенизационный контракт и справедливы независимо от модели.

Структура работы: раздел~\ref{sec:encoder} описывает энкодер и декодер;
раздел~\ref{sec:theory} доказывает две оценки; раздел~\ref{sec:loss}
вводит порядковую мягкую кросс-энтропию и мотивирует её через оценку
дифференциальной энтропии; раздел~\ref{sec:experiments} приводит
результаты на M4 Hourly и ETTh1.

\section{Токенизация на основе производных}
\label{sec:encoder}

Пусть $x \in \mathbb{R}^{T \times F}$ --- многомерный сигнал из $F$ каналов
длины $T$. Фиксируем порядок производной $k \ge 1$ (все результаты ниже
--- для $k = 1$); определяем $D_c = \max_{t} |x_{t+1,c} - x_{t,c}|$
поканально (\emph{область определения}) и поканальный масштаб
$s_c = (V-2)/(2 D_c)$, где $V \in 2\mathbb{Z}_{\ge 2}$ --- размер словаря.

\subsection{Энкодер.}
Для каждого канала энкодер вычисляет
\begin{align}
y_t &= s \cdot (x_{t+1} - x_t), \qquad |y_t| \le (V-2)/2, \\
u_t &= \mathrm{trunc}(y_t) \in \mathbb{Z}, \\
e_t &= y_t - u_t, \\
E_t &= \textstyle\sum_{s \le t} e_s, \quad R_t = \mathrm{round}(E_t), \quad R_0 = 0, \\
c_t &= R_t - R_{t-1}, \\
q_t &= u_t + c_t + V/2 \quad \in \{0, \dots, V-1\}.
\end{align}
Это $\Sigma\Delta$-модулятор первого порядка. Остаток от усечения $e_t$
накапливается; когда накопитель пересекает $\pm \tfrac12$, во
следующий токен вкладывается единичный перенос $c_t \in \{-1, 0, +1\}$.
Токены, выходящие за диапазон $[0, V-1]$ вследствие переноса, клиппируются;
замечание~\ref{rem:clip} обсуждает этот эффект.

\subsection{Декодер.}
По начальному значению $x_1$ и последовательности токенов $(q_t)$
декодер реконструирует
\begin{equation}
\hat x_{t+1} = \hat x_t + (q_t - V/2)/s, \qquad \hat x_1 = x_1.
\end{equation}

\section{Токенизационный контракт}
\label{sec:theory}

Два утверждения на уровне токенизации, оба не зависят от модели.

\subsection{Оценка реконструкции}

\begin{theorem}[реконструкция, $k=1$]
\label{thm:recon}
Для любого сигнала с $D = \max_t |x_{t+1} - x_t| > 0$ выполнено
\begin{equation}
\|\hat x - x\|_\infty \;\le\; \frac{D}{V-2} \;=\; \frac{\Delta}{2},
\quad\text{равномерно по } T.
\end{equation}
\end{theorem}

Доказательство опирается на два элементарных факта о $\Sigma\Delta$-цикле.

\begin{lemma}\label{lem:round}
Пусть $\rho_t = E_t - R_t$. Тогда $\rho_0 = 0$ и $|\rho_t| \le \tfrac{1}{2}$
для любого $t$.
\end{lemma}

\begin{proof}
Непосредственно: $R_t$ --- округление $E_t$ до ближайшего целого.
\end{proof}

\begin{lemma}[сокращение на шаге]\label{lem:step}
$u_t + c_t = y_t + \rho_{t-1} - \rho_t.$
\end{lemma}

\begin{proof}
$c_t = R_t - R_{t-1} = (E_t - \rho_t) - (E_{t-1} - \rho_{t-1}) = e_t + \rho_{t-1} - \rho_t$.
Прибавив к обеим сторонам $u_t$ и используя $y_t = u_t + e_t$, получаем требуемое.
\end{proof}

\begin{proof}[Доказательство теоремы~\ref{thm:recon}]
Декодированный прирост равен $\hat x_{t+1} - \hat x_t = (q_t - V/2)/s = (u_t + c_t)\Delta$.
Истинный прирост $x_{t+1} - x_t = y_t \Delta$. По лемме~\ref{lem:step} их
разность равна $(\rho_{t-1} - \rho_t)\Delta$. Суммируя от $t = 1$ до $T-1$,
получаем телескопическое сокращение:
\begin{equation}
\hat x_T - x_T \;=\; \Delta \sum_{t=1}^{T-1}(\rho_{t-1} - \rho_t)
             \;=\; \Delta\,(\rho_0 - \rho_{T-1})
             \;=\; -\rho_{T-1}\,\Delta,
\end{equation}
где использовано $\hat x_1 = x_1$ и $\rho_0 = 0$. По лемме~\ref{lem:round}
$|\rho_{T-1}| \le \tfrac12$. Применив тот же аргумент к произвольному
префиксу, получаем оценку, равномерную по $t$.
\end{proof}

\begin{corollary}[масштабирование]
$\|\hat x - x\|_\infty = O(D/V)$. Удвоение словаря уменьшает ошибку
реконструкции в два раза.
\end{corollary}

\begin{remark}[клиппирование]\label{rem:clip}
Поскольку $|y_t| \le (V-2)/2$ и $|c_t| \le 1$, имеем $q_t \in [0, V] \cup \{-1\}$.
Насыщение на границе $\{-1, V\}$ клиппируется и вносит не более одного
дополнительного $\Delta$ на клиппированный шаг. Так как $|y_t|$ достигает
максимума лишь на конечном (как правило --- пустом) множестве $t$, где ещё
и направление переноса совпадает, суммарный эффект имеет порядок
$O(\Delta)$ и поглощается в оценку скорости.
\end{remark}

\begin{remark}[старшие производные]
При $k \ge 2$ декодер применяет $k$ интегрирований. $\Sigma\Delta$
первого порядка компенсирует одно из них; оставшиеся $k-1$ превращают
ограниченный остаток в дрейф порядка $T^{k-1}\Delta/2$. Именно поэтому
текущий энкодер использует $k = 1$. $\Sigma\Delta$-модулятор порядка $k$
восстановил бы оценку $O(\Delta)$ ценой более сложного состояния.
\end{remark}

\subsection{Оценка кросс-энтропии}

Пусть $X_t$ --- непрерывный прирост на шаге $t$ при условии контекста
$\mathrm{ctx}$, с плотностью $p(\cdot \mid \mathrm{ctx})$ и дифференциальной
энтропией $h(X_t \mid \mathrm{ctx}) = -\int p \log p$.

\begin{proposition}[достижимая CE]
\label{prop:ce}
В пределе тонкого квантования $\Delta \to 0$ кросс-энтропия оптимального
авторегрессионного предиктора удовлетворяет
\begin{equation}
\mathcal{L}^\star(V) \;=\; h(X_t \mid \mathrm{ctx}) + \log V - \log(2D) \;+\; o(1).
\end{equation}
То есть $\mathcal L^\star$ линейна по $\log V$ с угловым коэффициентом $1$.
\end{proposition}

\begin{proof}[Набросок]
Стандартное соотношение между дискретной и дифференциальной энтропией:
при плотности $p$ и ширине ячейки $\Delta$
\begin{equation}
H(\mathrm{bin}) = -\sum_b P(\mathrm{bin}{=}b) \log P(\mathrm{bin}{=}b) = h(X) - \log \Delta + o(1)
\qquad (\Delta \to 0).
\end{equation}
(См., напр., Cover \& Thomas, \emph{Elements of Information Theory}, \S 8.3.)
Условная версия заменяет $X$ на $X \mid \mathrm{ctx}$. Подставляя
$\Delta = 2D/(V-2)$, получаем приведённую формулу.
\end{proof}

\begin{corollary}[инвариантная по словарю метрика мастерства]
\label{cor:skill}
Модели при $V_1 \ne V_2$ сопоставимы только через
\emph{остаток дифференциальной энтропии}
\begin{equation}
\tilde h(V) := \mathcal L(V) - \log V + \log(2D),
\end{equation}
который оценивает $h(X_t \mid \mathrm{ctx})$ и в пределе не зависит от $V$.
\end{corollary}

\subsection{Эмпирическая проверка}

На ETTh1 при seq\_len = 96, pred\_len = 96, 2000 итераций:

\begin{center}
\begin{tabular}{lccc}
\toprule
$V$ & val CE (nats) & $\log V$ & $\mathcal L(V) - \log V$ \\
\midrule
256 & 3{,}96 & $5{,}545$ & $-1{,}585$ \\
64  & 2{,}51 & $4{,}159$ & $-1{,}649$ \\
\bottomrule
\end{tabular}
\end{center}

Предсказанная разница $\log(256/64) = 1{,}386$ нат согласуется с
наблюдаемой $1{,}45$ нат с точностью до $0{,}06$ --- последнее представляет
собой остаток $\tilde h$. Это согласуется с предложением~\ref{prop:ce}.

\section{Порядковая мягкая кросс-энтропия}
\label{sec:loss}

Стандартная кросс-энтропия рассматривает $V$ бинов как номинальные
категории: прогноз, ошибающийся на один бин, штрафуется так же, как
прогноз, ошибающийся на пятьдесят. Это ложь: бины --- порядковый
индекс над непрерывной шкалой, и ``ошибка на один бин'' почти верна.

Заменяем one-hot цель $\mathbf 1_{i = y}$ на дискретную гауссиану
по индексам бинов:
\begin{equation}
p^{\text{soft}}_i(y; \sigma) \;\propto\; \exp\!\left(-\tfrac{(i - y)^2}{2\sigma^2}\right),
\qquad \sigma \ge 0,
\end{equation}
и обучаем стандартной мягкой кросс-энтропией
$\mathcal L_\sigma = -\sum_i p^{\text{soft}}_i \log q_i$, где
$q = \mathrm{softmax}(\text{logits})$. При $\sigma \to 0$ восстанавливается
обычная кросс-энтропия с one-hot целью.

\paragraph{Почему это важно для $\Sigma\Delta$-квантованного сигнала.}
Непрерывное значение, попавшее вблизи границы бина, с точностью до переноса
относится к одному из двух соседних бинов. $\Sigma\Delta$-энкодер может
сдвинуть его в любую сторону в зависимости от накопителя. Обучение с
one-hot целью делает вид, что это назначение детерминировано; мягкая
цель признаёт по существу стохастическое разрешение равенства. В
более общем случае любая алеаторная компонента $p(X \mid \mathrm{ctx})$
делает истинное условное распределение по бинам и так гладким,
похожим на гауссиану, и label smoothing с масштабом $\sigma$ \emph{в
единицах бина} --- это порядковый эквивалент подгонки под это распределение.

\paragraph{Бин-пространство против шкалы значений.}
Так как ширина бина $\Delta = 2D/V$, постоянное сглаживание в шкале
значений $\sigma_{\text{val}} = \sigma \cdot \Delta$ требует, чтобы
$\sigma$ росло \emph{линейно по} $V$. Величина $\sigma$, подобранная
при $V = 64$, окажется заниженной в $4$ раза при $V = 256$.

\section{Эксперименты}
\label{sec:experiments}

\subsection{M4 Hourly (краткосрочный, глобальный)}

Протокол M4: поканальная $z$-нормировка, одна глобальная модель,
обученная по всем $414$ Hourly-рядам, горизонт $H = 48$. Модель:
2-слойный GPT, $d = 64$, $h = 4$, $V = 256$, 20k итераций, AdamW,
восстановление best-val чекпоинта. Метрика: sMAPE (меньше --- лучше).

\begin{center}
\begin{tabular}{lcc}
\toprule
потеря & sMAPE & MASE \\
\midrule
обычная CE ($\sigma=0$)    & 18{,}58 & -- \\
мягкая CE $\sigma = 1.0$  & 17{,}55 & 1{,}663 \\
мягкая CE $\sigma = 2.0$  & \textbf{17{,}49} & \textbf{1{,}659} \\
мягкая CE $\sigma = 3.0$  & 17{,}35 & 1{,}949 \\
\bottomrule
\end{tabular}
\end{center}

$\sigma = 2$ --- совместный оптимум по обеим метрикам. При $\sigma = 3$
sMAPE продолжает улучшаться, но MASE резко ухудшается --- классический
признак пере-сглаживания, когда прогнозы стягиваются к условному среднему,
что выгодно для симметричного sMAPE, но проигрывает под наивно-сезонной
нормировкой MASE. Выигрыш в один пункт sMAPE при $\sigma = 1.0$ по
сравнению с обычной CE --- ключевой качественный вывод: \emph{на
короткогоризонтном бенчмарке с множеством коротких разнородных рядов
сказать модели ``ошибиться на один бин почти не ошибка'' стоит столько же,
сколько заметное архитектурное изменение.}

\subsection{ETTh1 (долгосрочный, многоканальный)}

Протокол iTransformer: $7$ каналов, seq\_len = $96$, pred\_len = $96$,
хронологический раздел $12/4/4$ месяца, StandardScaler по train.
Одна модель при $d = 128$, $h = 4$, $2$ слоя, $2000$ итераций, batch $32$, AdamW.

\begin{center}
\begin{tabular}{lcc}
\toprule
конфигурация & MSE & MAE \\
\midrule
$V = 256$, обычная CE                 & 0{,}9861 & 0{,}6056 \\
$V = 64$, обычная CE                   & 0{,}9787 & 0{,}6030 \\
$V = 64$, мягкая CE $\sigma = 1.0$     & 0{,}9806 & 0{,}6034 \\
\bottomrule
\end{tabular}
\end{center}

ETTh1 упирается в потолок по ёмкости / подгонке на уровне
$\mathrm{MSE} \approx 0{,}98$: и уменьшение словаря, и мягкие цели
двигают результат в пределах шума. Контраст с M4 наводит на мысль,
что \emph{ценность} мягких целей сосредоточена в режиме нехватки данных,
когда множество коротких рядов вынуждает модель экстраполировать
условные распределения без достаточного сигнала, чтобы локализовать
конкретный бин.

\subsection{Сводка по масштабированию словаря и функции потерь}

Бюджет токенизации против бюджета потерь:
\begin{center}
\begin{tabular}{lcc}
\toprule
                         & на удвоение словаря & режим \\
\midrule
идеальная CE $\mathcal L^\star(V)$   & $+\log 2 \approx 0{,}69$ нат & предл.~\ref{prop:ce} \\
реконструкция $\|\hat x - x\|_\infty$ & $\times \tfrac12$ & теор.~\ref{thm:recon} \\
порядковое сглаживание $\sigma^\star$ & $\times 2$ (сохранение в шкале значений) & \S\ref{sec:loss} \\
\bottomrule
\end{tabular}
\end{center}

\section{Связанные работы}

\textbf{Инвертированные трансформеры и channel-mixing.} iTransformer
\cite{itransformer2024} и PatchTST \cite{patchtst2023} отказываются от
авторегрессионной головы и регрессируют весь горизонт за один проход ---
доминирующая схема на долгогоризонтных бенчмарках. Наша модель
оставляет рецепт GPT неизменным; индуктивное смещение переносится
на токенизацию.

\textbf{Порядковая классификация и label smoothing.} Label smoothing
\cite{szegedy2016} равномерно распределяет вероятность по нецелевым
классам. Варианты со взвешиванием по расстоянию (cumulative link models,
ordinal softmax) встречаются в литературе по порядковой регрессии.
Гауссовские мягкие цели появляются, в частности, в CORAL \cite{cao2020}
и в score-matching диффузионных моделей для дискретных целей
\cite{austin2021}. Наш вклад --- привязка масштаба сглаживания $\sigma$
к токенизации на основе $\Sigma\Delta$-энкодера.

\textbf{$\Sigma\Delta$-квантование.} $\Sigma\Delta$-модуляторы первого
порядка --- классический элемент АЦП \cite{candy1992}; мы применяем их
как недифференцируемый шаг предобработки перед трансформером.

\section{Заключение}

Decoder-only трансформер, авторегрессионно обученный на токенах
производной с $\Sigma\Delta$-квантованием, достигает конкурентных
результатов в краткосрочном и долгосрочном прогнозировании без
специализированной головы. Токенизационный контракт даёт два
замкнутых закона масштабирования словаря; порядковая мягкая
кросс-энтропия согласована с естественной геометрией по бинам и
приносит относительное улучшение sMAPE на $6\%$ на M4 Hourly
без изменения архитектуры.

\begin{thebibliography}{9}
\bibitem{itransformer2024} Liu, Y.\ et al.\ \emph{iTransformer: Inverted
Transformers are Effective for Time Series Forecasting}. ICLR 2024.
arXiv:2310.06625.

\bibitem{patchtst2023} Nie, Y.\ et al.\ \emph{A Time Series is Worth 64 Words:
Long-term Forecasting with Transformers}. ICLR 2023.

\bibitem{szegedy2016} Szegedy, C.\ et al.\ \emph{Rethinking the Inception
Architecture for Computer Vision}. CVPR 2016.

\bibitem{cao2020} Cao, W., Mirjalili, V., Raschka, S.\ \emph{Rank-consistent
Ordinal Regression for Neural Networks with Application to Age Estimation}.
Pattern Recognition Letters 2020.

\bibitem{austin2021} Austin, J.\ et al.\ \emph{Structured Denoising Diffusion
Models in Discrete State-Spaces}. NeurIPS 2021.

\bibitem{candy1992} Candy, J.C., Temes, G.C.\ \emph{Oversampling Delta-Sigma
Data Converters}. IEEE Press 1992.
\end{thebibliography}

\end{document}
